% SHARD-ID: AQM-86-QSA-DEGENERATE-OVERRESTRICTIVE-CASES
% SHARD-TITLE: Degenerate and Over-Restrictive Cases
% SHARD-SUMMARY: Catalogues concrete QSA failure modes without promoting them to general theorems.
% SHARD-SUMMARY: Separates collapsed quotients, radical quotients, missing choices, zero outputs, empty inputs, blocked rows, and not-applicable workflows.
% SHARD-SUMMARY: Carries the Step 21 blockers forward and keeps degenerate cases visible in the catalogue.
% SHARD-KEYWORDS: quantum-system association, degeneracy, over-restrictive, radical, quotient collapse, zero, empty, blocked

\section{Degenerate and Over-Restrictive Cases}
\label{sec:qsa-degenerate-overrestrictive-cases}

\subsection{Status and Local Anchors}

\textbf{Status: New failure-mode catalogue.}  This is QSA Step 22 from
\path{docs/quantum_system_associations/PLAN.md}.  It adds no source,
convention entry, script, run bundle, generated data, populated database row,
Beads update, or report-level classification theorem.  It records concrete
failure modes already visible in AQM-65--AQM-85.

The vocabulary and evidence rule are AQM-65 and \texttt{CONVENTIONS.md}
convention \texttt{(ap)}.  The status tokens are exactly those of convention
\texttt{(au)}.  The main anchors are AQM-67--AQM-71 for finite-set and
radical examples, AQM-74 for single-particle comparison boundaries, AQM-77
for cotangent zero and cotangent-choice examples, AQM-80--AQM-83 for
finite-ring, quotient, nilpotent, and product-ring rows, AQM-84 for infinite
boundary statuses, and AQM-85 for the gap-first agreement ledger.

\subsection{Reading Rule}

The word \texttt{degenerate} is not a dismissal.  It marks a concrete row in
which a workflow collapses, forgets structure, becomes too small for an
intended comparison, or reduces to an earlier row.  Such rows stay in the
catalogue because they are useful negative controls.

The token \texttt{blocked} is different.  It means that required evidence,
conventions, or tooling are missing.  A blocked row is not a mathematical
obstruction and is not evidence that the workflow fails.  Likewise,
\texttt{unknown} means not locally determined, and \texttt{not\_applicable}
means that the workflow is outside the row's stated semantics.

\subsection{Too Little Structure}

\paragraph{Bare finite set without a Gram choice.}
For a two-point set \(X=\{a,b\}\), AQM-67 and convention \texttt{(aq)}
produce only
\[
  \mathbb C\langle X\rangle=\mathbb C e_a\oplus\mathbb C e_b
\]
with formal basis labels.  No inner product, adjoint, observable algebra,
Hamiltonian, tensor factorization, or Weyl label space is returned.  The
identity Gram matrix would be extra data.  Therefore any Hilbert-space or
unitary comparison from the bare set alone is \texttt{blocked} by missing
Gram data, not proved false.

\paragraph{Empty input versus zero output.}
AQM-67 gives
\[
  \mathbb C\langle\emptyset\rangle=0 .
\]
Here \(\emptyset\) is the \texttt{empty} input and \(0\) is the \texttt{zero}
carrier output.  AQM-69 gives the same zero carrier for the empty direct sum,
whereas AQM-68 gives the empty tensor product \(\mathbb C\).  AQM-74 uses
this as a concrete obstruction to silently identifying tensor-site
coexistence with direct-sum alternatives.

\subsection{Radicals and Quotient Label Spaces}

\paragraph{Affine maps on \(\mathbb F_3\).}
AQM-71 reviews the checked row displayed below.
\[
  \texttt{vector\_space\_F3\_affine}
\]
The row is in
\path{runs/2026-05-24-arithmetic-quantum-fields/data/arithmetic_quantum_field_examples.csv}.  The scalar space has
basis \(1,t\), the field phase dimension is \(4\), the symplectic radical has
dimension \(2\), and the reduced phase dimension is \(2\).  The honest Weyl
label space is therefore the quotient \(E/\operatorname{rad}(E)\), not the
unreduced affine-function label space.  Forgetting this quotient would
overcount labels.

\paragraph{Linear coordinate functions on \(\mathbb F_3^2\).}
The same AQM-71 review lists the checked row displayed below.
\[
  \texttt{affine\_plane\_F3\_linear}
\]
It has field phase dimension \(4\),
symplectic rank \(0\), radical dimension \(4\), and reduced phase dimension
\(0\).  This is a concrete \texttt{degenerate} radical case: the selected
linear coordinate-function space is totally radical for the all-weights-one
pairing, so the reduced Weyl label group is trivial and the Hilbert carrier
is one-dimensional.  The row is not removed from the catalogue; it records
that the chosen field space and pairing were too restrictive for nontrivial
Weyl labels.

\paragraph{Projective-line sheaf sections.}
AQM-16 gives a second checked radical warning.  For
\(\mathbb P^1_{\mathbb F_3}\) with \(\mathcal O(1)\), the rational-point
evaluation row has scalar pairing rank \(0\), symplectic rank \(0\), radical
dimension \(4\), and reduced phase dimension \(0\) in
\path{runs/2026-05-24-projective-line-sheaf-fields/data/projective_line_sheaf_field_summary.csv}.
For \(\mathcal O(4)\), the evaluation kernel and radical are nonzero but the
reduced phase dimension is still \(8\).  These examples show why the radical
quotient is part of the workflow, not cleanup after the fact.

\subsection{Zero Cotangent Labels}

\paragraph{The field \(\mathbb F_3\).}
AQM-77 records that for \(X=\Spec\mathbb F_3\) over
\(\mathbb F_3\), both the intrinsic cotangent and the relative cotangent are
zero:
\[
  C_x^{\rm loc}=C_x^{\rm rel}=0 .
\]
The cotangent-first phase label is therefore zero and the Weyl carrier is
one-dimensional.  This is a \texttt{zero} output for that workflow.  It is
not the residue-site workflow, where AQM-81 has one qutrit site.

\paragraph{Product fields.}
For \(R=k^n\) with \(k=\mathbb F_3\), AQM-83 cites AQM-58 and AQM-77:
\[
  \Omega^1_{k^n/k}=0
\]
and the local cotangent labels vanish at every product point.  The same ring
has \(n\) residue-field Weyl factors in AQM-40 and AQM-83.  Thus the
cotangent-first row is \texttt{degenerate}/\texttt{zero} for the selected
cotangent labels, while the residue-site row remains nontrivial.  The
catalogue must keep both rows.

\subsection{Quotients and Collapsed Presentations}

\paragraph{\(\mathbb Z/6\mathbb Z\) versus the product presentation.}
AQM-23 computes the two residue fields \(\mathbb F_2\) and
\(\mathbb F_3\) for \(\Spec(\mathbb Z/6\mathbb Z)\), and AQM-81/AQM-83 use
the mixed-residue carrier
\[
  \ell^2(\mathbb F_2)\otimes\ell^2(\mathbb F_3).
\]
AQM-80 records that the finite-ring review export treats the
\(\mathbb F_2\times\mathbb F_3\) presentation as merged to the
\(\mathbb Z/6\mathbb Z\) representative.  This is a concrete
\texttt{degenerate} catalogue warning: the two presentations should not be
counted as independent examples in this local evidence layer.  It is not a
general quotient-collapse theorem.

\paragraph{The \(\mathbb F_2[x]/(x^2+x)\) candidate.}
AQM-80 lists \(\mathbb F_2[x]/(x^2+x)\) only as a possible
quotient-collapse row.  It remains \texttt{blocked} by
\texttt{aqm-qsa-collapse01}; the current report has no local derivation or
certificate that may be promoted to a checked collapse row.  This is a
concrete unsupported finite-ring example, and its blocked status is part of
the catalogue.

\subsection{Nilpotents Forgotten by Reduced Layers}

\paragraph{Dual numbers.}
For
\[
  A=\mathbb F_3[\epsilon]/(\epsilon^2),
\]
AQM-82 records one reduced residue site with carrier \(\ell^2(\mathbb F_3)\).
The nilpotent class \(\epsilon\) reduces to zero at that site.  AQM-36 and
AQM-82 also record the thickened Artin carrier
\[
  \ell^2(\mathbb F_3[\epsilon]/(\epsilon^2)),
  \qquad \dim_{\mathbb C}=9,
\]
after the top-coefficient generating character is fixed.  The reduced row is
\texttt{degenerate} if the intended question is nilpotent-sensitive; the
thickened row is a separate available local Artin workflow.

\paragraph{The \(t^n=1\) family.}
AQM-82 reviews AQM-35/AQM-37 for
\[
  R_n=\mathbb F_3[t]/(t^n-1),\qquad n=3^s m,\quad 3\nmid m .
\]
The reduced residue layer has Hilbert dimension \(3^m\), while the locally
proved thickened layer has dimension \(3^n\).  The repeated factors are not
extra closed points; they are thickened degrees at the existing reduced
points.  Removing this row would hide exactly the nilpotent-sensitive
failure mode Step 22 is meant to catalogue.

\subsection{Mixed Residues and Not-Applicable Readings}

The \(\mathbb Z/6\mathbb Z\) row is also a concrete
\texttt{not\_applicable} warning.  AQM-23, AQM-80, AQM-81, and AQM-83 use
the mixed label group
\[
  \mathbb F_2^2\oplus\mathbb F_3^2 .
\]
A single \(\mathbb F_3\)-linear field-label reading is outside the row's
semantics.  This is not a blocked theorem about impossibility; it is the
wrong workflow specification for the selected mixed-residue object.

Similarly, AQM-82 marks the nilpotent-sensitive workflow as
\texttt{not\_applicable} for the field row \(\mathbb F_3\): there is no
nilpotent thickening in that object.  The status is different from the
blocked thickened rows for unsupported finite rings.

\subsection{Unsupported Infinite-Representation Cases}

AQM-84 gives the concrete infinite boundary rows
\(\Spec\mathbb F_q[t]\), \(\Spec\mathbb F_q[x,y]\), and \(\Spec\mathbb Z\).
Finite closed supports and algebraic quasi-local closed-point algebras are
available there.  The following are not available finite-site conclusions:

\begin{itemize}
\item generic and curve-generic residue fields as finite Weyl sites;
\item local-field, adelic, profinite, analytic, or \(C^*\)-completed
  representations without the missing topology, character, norm, and
  representation conventions;
\item formal infinite-volume regular-function profiles as actual
  quasi-local Weyl operators.
\end{itemize}

These are \texttt{blocked}, generic-sector proposal, formal infinite-volume,
or completion-proposal statuses in the AQM-84 vocabulary.  They are not
mathematical no-go results.

\subsection{Status-Token Checklist}

The concrete distinctions to preserve in later shards are:

\begin{description}
\item[\texttt{empty}.] The indexing input or support is empty, as in
  \(X=\emptyset\) for AQM-67 or the empty finite support in a local net.
\item[\texttt{zero}.] The mathematical output is the zero carrier or zero
  label object, as in \(\mathbb C\langle\emptyset\rangle=0\) or the zero
  cotangent labels of \(\Spec\mathbb F_3\).
\item[\texttt{degenerate}.] A concrete row collapses or forgets structure:
  total radical reduction, product-field zero cotangent labels, reduced
  residue layers forgetting nilpotents, or merged \(\mathbb Z/6\mathbb Z\)
  product presentations.
\item[\texttt{blocked}.] Evidence, conventions, or tooling are missing, as
  for \(\mathbb F_2[x]/(x^2+x)\) quotient collapse, arbitrary finite-ring
  residue decomposition, database-backed cotangent rows, and infinite
  completions.
\item[\texttt{unknown}.] No local determination has been recorded, as in
  AQM-80's missing cotangent comparison row for \(\mathbb Z/6\mathbb Z\).
\item[\texttt{not\_applicable}.] The workflow is outside the row semantics,
  as for a common \(\mathbb F_3\)-linear reading of the mixed-residue
  \(\mathbb Z/6\mathbb Z\) row or a nilpotent-sensitive thickening of the
  field \(\mathbb F_3\).
\end{description}

\subsection{Blockers Carried Forward}

The open blockers from AQM-85 remain unchanged:
\[
  \begin{gathered}
  \texttt{aqm-qsa-src02},\quad
  \texttt{aqm-qsa-frres01},\quad
  \texttt{aqm-qsa-frcot01},\\
  \texttt{aqm-qsa-collapse01},\quad
  \texttt{aqm-qsa-artin01}.
  \end{gathered}
\]
This shard does not close or reinterpret any of them.  Step 23 should turn
these preserved failure modes into explicit open problems and next catalogue
targets, not into unsupported equivalence or no-go theorems.
